\section{Introduction}\label{sec:introduction}

% SOURCE: README.md and research_outputs/hodgecy_ii/final/hodgecy_ii_final_results.*
% TARGET LENGTH: 4--5 pages.
% CLAIM: HodgeCY II studies source-fidelity and structural block-evaluation
% collapse in the frozen v1.0.0 double-octic cohort.

\subsection{The Problem}

Introduce double-octic presentations as a setting where coarse invariants can
collapse distinct source data.  Emphasize that Hodge-number degeneracy is
classical and not the novelty claim.

\subsection{Main Contributions}

\begin{theorem}[Source-fidelity census]\label{thm:source-fidelity-census}
For the frozen HodgeCY v1.0.0 double-octic cohort, the HodgeCY II source
census processes \(456\) presentations and records \(114\) nontrivial
source-level pairs or larger sets.  These decompose as \(57\) pairs,
\(13\) triples, and \(44\) larger sets.
\end{theorem}

\begin{proof}[Evidence]
This is a frozen computational census statement.  The reproducible evidence is
the final HodgeCY II synthesis bundle and the deterministic reproduction script
`scripts/reproduce_hodgecy_ii.py`.
\end{proof}

\begin{proposition}[The \(84/84a\) source witness]\label{prop:source-witness-84-84a}
The presentations \(84\) and \(84a\) have the same local inventory, the same
Hodge data \((h^{1,1},h^{1,2})=(40,0)\), and the same rational source rank
\(26\).  They are separated by integral source type:
\[
  {\rm SNF}(84)=(1^{23},2,6,12),
  \qquad
  {\rm SNF}(84a)=(1^{21},2,4,4,4,12).
\]
\end{proposition}

\begin{proof}[Evidence]
This is the source-lattice comparison certified in the HodgeCY II final
evidence bundle.  The proposition is a statement about the verified source
assembly records only.
\end{proof}

\begin{theorem}[Eight-plane line skeleton]\label{thm:line-skeleton}
Let \(S=k[x_0,x_1,x_2,x_3]\), and let
\(L_1,\ldots,L_8\in S_1\) be distinct linear forms such that every pair has
rank \(2\) and every triple has rank \(3\).  Equivalently, no three
corresponding planes contain a common line.  If \(F=\prod_i L_i\) and
\[
  I_C=(F/L_1,\ldots,F/L_8),
\]
then
\[
  I_C=\bigcap_{1\le i<j\le 8}(L_i,L_j).
\]
Thus \(I_C\) is the saturated reduced ideal of the union of the \(28\)
pairwise intersection lines, and it has Hilbert--Burch resolution
\[
0\longrightarrow S(-8)^7\longrightarrow S(-7)^8
\longrightarrow I_C\longrightarrow 0 .
\]
Four-plane, or higher, point concurrences are allowed: the argument only uses
the codimension-two condition that no three planes contain a common line.
\end{theorem}

\begin{proof}[Proof source]
Use the self-contained induction in
`research_outputs/hodgecy_ii/final/hodgecy_ii_hilbert_burch_theorem.tex`.
For the manuscript body, reproduce that proof rather than replacing it by an
appeal to the full star-configuration theorem of
Geramita--Harbourne--Migliore~\cite{GeramitaHarbourneMigliore2013}.
\end{proof}

\begin{proposition}[Regular quartic block section and mapping cone]\label{prop:quartic-block-section}
For the verified HodgeCY II arrangements \(84\) and \(84a\), the frozen quartic
block section is regular on the eight-plane line skeleton.  If \(I_B=I_C+(Q)\),
then multiplication by \(Q\) gives an exact sequence
\[
0\longrightarrow (S/I_C)(-4)\longrightarrow S/I_C
\longrightarrow S/I_B\longrightarrow 0,
\]
and \(S/I_B\) has mapping-cone resolution
\[
0
\longrightarrow S(-12)^7
\longrightarrow S(-11)^8 \oplus S(-8)^7
\longrightarrow S(-7)^8 \oplus S(-4)
\longrightarrow S
\longrightarrow S/I_B
\longrightarrow 0 .
\]
\end{proposition}

\begin{proof}[Proof source]
This is the regular quartic block-section theorem and mapping-cone corollary
in the final Hilbert--Burch theorem asset.  The verifier checks nonzero
squarefree restriction of \(Q\) on all \(28\) associated line primes for both
frozen arrangements.
\end{proof}

\begin{corollary}[Verified block Hilbert series]\label{cor:block-hilbert-series}
For the verified HodgeCY II block schemes attached to \(84\) and \(84a\),
\[
  {\rm Hilb}_{S/I_B}(t)=
  \frac{(1-t^4)(1-8t^7+7t^8)}{(1-t)^4}.
\]
Consequently
\[
H_B(d)=1,4,10,20,34,52,74,92,105,112,112,\ldots,
\]
with stabilization at degree \(9\).
\end{corollary}

\begin{proof}[Proof source]
This follows from the verified regular quartic block section and the mapping
cone argument in the final Hilbert--Burch theorem asset.
\end{proof}

\begin{corollary}[Degree-eight block-evaluation deficiency]\label{cor:block-deficiency}
For the verified HodgeCY II block schemes attached to \(84\) and \(84a\),
\[
  h^1({\bf P}^3,\mathcal I_B(8))=7.
\]
Equivalently, the degree-eight block-evaluation cokernel has dimension \(7\).
This is a block-scheme deficiency statement only.
\end{corollary}

\begin{proposition}[Source versus block-evaluation non-determination]\label{prop:source-block-nondetermination}
On the witness pair \(84/84a\), the verified degree-eight block-evaluation data
do not determine the integral source-assembly type.  The source layer splits by
Smith type, while the verified block Hilbert/evaluation data agree through
degree \(8\).
\end{proposition}

\begin{proof}[Evidence]
Combine Proposition~\ref{prop:source-witness-84-84a} with
Corollaries~\ref{cor:block-hilbert-series} and~\ref{cor:block-deficiency}.
This does not assert a reverse implication, a canonical source-to-evaluation
map, an integral evaluation lattice, an LMHS identification, or a classical
nodal-defect theorem.
\end{proof}


\subsection{Claims Firewall}

Summarize the firewall early and point to Section~\ref{sec:scope-nonclaims}.

\subsection{The HodgeCY Viewpoint}

Explain the source-aware approach: records, invariants, certificates, and the
separation between data layers.

\subsection{Relation to HodgeCY I}

State the continuity compactly.  HodgeCY I introduced local profiles/source
assembly and identified the \(84/84a\) rational-collapse/integral-separation
phenomenon; HodgeCY II places that witness in a broader census and follows it
into verified block geometry.

\subsection{Relation to Prior Work}

Preview the related-work firewall: Hodge-number degeneracy is classical, while
the source-fidelity hierarchy is the present contribution.

\subsection{Organization}

List the remaining sections and their roles.

\section{Background and Fidelity Hierarchy}\label{sec:background}

% SOURCE: docs/hodgecy_ii_cohort.md, docs/source_to_node_comparison.md,
% research_outputs/hodgecy_ii/manuscript_assets/figures/fidelity_hierarchy.svg.
% TARGET LENGTH: 4--5 pages.

\subsection{Double-Octic Arrangement Records}

Define the hierarchy used in the paper: local inventory, Hodge tuple, rational
source assembly, integral Smith type, torsion type, equivariant type, and later
block/evaluation layers.  Keep source-level invariants distinct from geometric
node, vanishing-cycle, LMHS, and Hodge-atom relations.

\subsection{Local Singularity Inventory}

Describe local inventory as the first coarse source-level comparison.

\subsection{Hodge Signatures}

Record the Hodge tuple layer and its role as a coarse collapse invariant.

\subsection{Rational Source Assembly}

Define the rational source-assembly comparison used for the \(84/84a\)
witness.

\subsection{Integral Source Assembly and Smith Type}

Introduce integral Smith type and torsion factors as the first separating
layer for \(84/84a\).

\subsection{Equivariant Source Layer}

Describe the equivariant layer as a finer source-level comparison without
turning it into a Hodge-theoretic relation.

\subsection{The Fidelity Hierarchy}

Use Figure II.1 to summarize the ladder.

\section{The Frozen v1.0.0 Double-Octic Cohort}\label{sec:cohort}

% SOURCE: research_outputs/hodgecy_ii/census/* and
% research_outputs/hodgecy_ii/complete_fidelity_pairs_and_sets.*
% TARGET LENGTH: 4--5 pages.
% TABLE: II.1 Fidelity census summary.
% TABLE: II.2 Representative fidelity controls.
% FIGURE: II.1 Fidelity hierarchy.

State the cohort size \(456\), the number of nontrivial pairs/sets \(114\), and
the split \(57/13/44\).  Explain what was processed and what "nontrivial"
means.  Do not describe the 114 sets as geometrically classified.

\subsection{The 456-Presentation Census}

Define the frozen cohort and evidence sources.

\subsection{The 114 Nontrivial Fidelity Sets}

Summarize pairs, triples, and larger sets.

\subsection{First-Separation Statistics}

Describe where coarse agreement first breaks under finer invariants.

\subsection{Representative Controls}

Use \(61/451\), \(84/84a\), \(452/453\), \(84/240\), \(84a/239\), and
\(239/240/241\).  Preserve the \(451\) factor-normalization warning and the
\(452/453\) quadratic-field exact-validation deferral.

\subsection{Selection of \(84/84a\) as the Deep Witness}

Explain why the paper follows this pair into exact block geometry.

\section{The \(84/84a\) Source Witness}\label{sec:source-witness}

% SOURCE: final source-lattice comparison and manuscript table II.6.
% TARGET LENGTH: 4--5 pages.
% TABLE: II.3 84-neighborhood refinement.
% FIGURE: II.2 84-neighborhood refinement tree.

Present \(84/84a\) as the principal witness: same local inventory, same Hodge
data \((40,0)\), same rational source rank \(26\), but different integral Smith
types \((1^{23},2,6,12)\) and \((1^{21},2,4,4,4,12)\).

\subsection{Common Local Inventory}

\subsection{Common Hodge Data}

\subsection{Rational Source Collapse}

\subsection{Integral Smith Separation}

\subsection{The \(239/240/241\) Neighborhood}

\subsection{What Source Equality Does and Does Not Imply}

\section{The Eight-Plane Line Skeleton}\label{sec:line-skeleton}

% SOURCE: docs/verified_node_block_geometry.md and final block-geometry
% evidence.
% TARGET LENGTH: 4--5 pages.
% FIGURE: II.1 Fidelity hierarchy, recalled only if notation needs a diagram.

Set up the eight-plane skeleton and prove the direct codimension-two ideal
identity that underlies the structural theorem.

\subsection{Arrangement Setup}

\subsection{Pair and Triple Rank Hypotheses}

\subsection{Higher Point Concurrence}

\subsection{Relation to Star Configurations}

\subsection{The Direct Line-Ideal Identity}

\subsection{The Explicit Hilbert--Burch Matrix}

\subsection{Minimal Free Resolution}

\subsection{Hilbert Series of the Skeleton}

\section{Quartic Block Geometry}\label{sec:block-geometry}

% SOURCE: docs/verified_node_block_geometry.md,
% docs/block_scheme_evaluation.md, and
% research_outputs/hodgecy_ii/final/hodgecy_ii_hilbert_burch_theorem.tex.
% TARGET LENGTH: 6--7 pages.
% FIGURE: II.4 Hilbert-profile comparison.
% TABLE: II.5 84/84a Hilbert/evaluation comparison.
% TABLE: II.4 84/84a block/node certification status.
% FIGURE: II.3 Block/node certification bridge.

Describe the verified degree-112 block schemes attached to the two
presentations.  State that the ordinary-node promotion and full saturated
Jacobian-node equality remain open.

\subsection{The Frozen Quartic Perturbation}

\subsection{Restriction to the 28 Lines}

\subsection{Squarefreeness and Regularity}

\subsection{The Verified Degree-112 Block Scheme}

\subsection{Mapping-Cone Resolution}

\subsection{Structural Hilbert Series}

\subsection{Stabilization at Degree 9}

\section{Critical-Degree Evaluation}\label{sec:critical-evaluation}

% SOURCE: final block-evaluation evidence and docs/block_scheme_evaluation.md.
% TARGET LENGTH: 3--4 pages.
% FIGURE: II.5, supplementary unless needed in the main text.
% TABLE: II.5 84/84a Hilbert/evaluation comparison.

\subsection{Degree-Eight Ambient Section Space}

\subsection{Hilbert Value \(H_B(8)=105\)}

\subsection{Evaluation Rank and Cokernel}

\subsection{Cohomological Interpretation}

\subsection{\(h^1(\mathcal I_B(8))=7\)}

\subsection{Structural Origin of the Number Seven}

\subsection{Classical Defect Firewall}

\section{Source Versus Block-Evaluation}\label{sec:source-block}

% SOURCE: final source-block comparison evidence and table II.6.
% TARGET LENGTH: 3--4 pages.
% TABLE: II.6 84/84a source versus block-evaluation comparison.
% FIGURE: II.6 Source versus block-evaluation axes.

Combine the integral source split with the identical verified block
Hilbert/evaluation data.  State the non-determination proposition exactly and
do not infer a source-to-evaluation morphism.

\subsection{The Two-Axis Comparison}

\subsection{The Non-Determination Witness}

\subsection{What the Structural Collapse Forgets}

\subsection{Rational Versus Integral Information}

\subsection{No Source-to-Evaluation Map Is Asserted}

\subsection{Integral Evaluation Lattice as an Open Problem}

\section{Related Work}\label{sec:related-work}

% SOURCE: related_work_notes.md and references.bib.
% TARGET LENGTH: 3--4 pages.

Position the work against known large Hodge-number degeneracy, double-octic
period refinements, degeneration-sensitive invariants, topological refinement
of Kreuzer--Skarke data, and Hilbert--Burch/star-configuration context.

\subsection{Kreuzer--Skarke and Large Hodge Degeneracy}

\subsection{Statistical Degeneracy in Calabi--Yau Data}

\subsection{Double-Octic Period Refinements}

\subsection{Degeneration-Sensitive Hodge-Theoretic Invariants}

\subsection{Cross-Construction Topological Refinements}

\subsection{Hilbert--Burch and Star-Configuration Context}

\section{Status Matrix and Future Work}\label{sec:status-future}

% SOURCE: final evidence/claim-status matrix and open problems.
% TARGET LENGTH: 2--3 pages.
% TABLE: II.7 Final evidence/claim-status matrix.

State what is proved, what is computationally certified, what is conditional,
and what remains open.  Route ordinary-node promotion, classical nodal defect,
source-to-node comparison, integral evaluation lattices, LMHS/MHM/Hodge atoms,
and full population geometry to future work.

\subsection{Evidence Hierarchy}

\subsection{Exact Arithmetic}

\subsection{Theorem Certificates}

\subsection{Reproduction Workflow}

\subsection{Claim-Status Matrix}

\section{Discussion}\label{sec:discussion}

% TARGET LENGTH: 2--3 pages.

\subsection{What HodgeCY II Establishes}

\subsection{Source-Aware Calabi--Yau Classification}

\subsection{Relation to Conifold and Hodge-Atom Questions}

\subsection{Open Geometric Promotion Problems}

\subsection{HodgeCY III Handoff}

\section{Scope and Nonclaims}\label{sec:scope-nonclaims}

The HodgeCY II claims are deliberately narrower than the surrounding
geometric questions.  The verified statements are source-fidelity census
results, the \(84/84a\) source witness, exact degree-112 block-scheme
certification, the Hilbert--Burch structural explanation for the verified
block schemes, and the resulting source-versus-block-evaluation
non-determination statement.

The manuscript does not claim any of the following:

\begin{enumerate}
\item a complete geometric analysis of all \(114\) nontrivial source-level
  pairs or sets;
\item ordinary-node promotion for the \(84/84a\) block schemes;
\item equality between the verified block schemes and full saturated Jacobian
  node schemes;
\item an unconditional classical nodal-defect theorem;
\item a natural source-to-evaluation or source-to-node chain map;
\item an integral evaluation-lattice theorem;
\item a source-to-vanishing-cycle map;
\item an LMHS, MHM, or Hodge-atom identification;
\item a reverse implication from block-evaluation data to source assembly;
\item projective equivalence or inequivalence solely from matching or differing
  computed signatures.
\end{enumerate}

The status of each item should be read as open unless the final evidence matrix
states otherwise.  In particular, the equality \(h^1({\bf P}^3,\mathcal
I_B(8))=7\) is a theorem for the verified block schemes, not a promotion of
the classical nodal defect.

